We asked a narrow question with instruments built so that the easy failure modes do not pay: a semantics of justified belief that keeps alternatives open, a contract that scores behaviour rather than representation, and a benchmark in which the world, the reports and what the reports justify are separable. The results are these. Retraction propagation through justifications, not ordinary belief updating, prevents downstream errors, by about twenty points in every setting including real revision histories. An incremental store that pins justifications returns the same status, value and alternatives as a symbolic replay does, with lower measured query latency, and saves total work only when queries are at least as frequent as observations or histories are long; its write cost grows exponentially in reports per key. Symbolic maintenance beats a prompted LLM adjudicator over the same extracted claims by a wide margin, but a rule-based adjudicator over the same claims does as well, so that result is about the adjudicators, not about write time. Refusing to assume recency buys protection exactly where the latest report is wrong, and costs abstention where it is right, which is most of the time in our generator and all of the time in a self-report benchmark; a pre-registered same-origin self-update rule closes the single-source case at a measured cost of at most 1.3 points of wrong commitments where the latest report is wrong. The evidence model's support is a defined, auditable quantity and not a better-calibrated or a better-committing one than simpler alternatives, and when the latest report is an injection the store commits to it on \rgPoisPalInj{} of queries, little less than recency. Everything needed to reproduce, contest or extend these results is released with the dated pre-registration and the cost ledger.
